\section{Open questions}

\subsection*{Q1. Does the headline survive on the real Kaggle file at natural imbalance?}
\textbf{Basis.} This replication used a synthetic BankSim-\emph{like} 200-record balanced
slice because Kaggle API access was unavailable in the subagent context. The paper
itself uses a hand-balanced 100/100 subset that erases the real-world $\sim 1.2\%$
fraud rate. On imbalanced data, F1 collapses and PR-AUC becomes the honest metric;
neither paper nor this replication reports PR-AUC.\\
\textbf{Next steps.} Obtain a Kaggle token, pull \texttt{ealaxi/banksim1}
($\sim$600K rows), rerun the 6-cell QSVC/VQC grid at (a) the paper's balanced protocol
and (b) natural imbalance with stratified sampling, reporting PR-AUC and precision at
recall = 0.90. Add XGBoost + Random Forest baselines. Budget: $\sim$2 h CPU (QSVC kernel
grows quadratically in $n_{\text{train}}$; subsample to 2000 if needed).

\subsection*{Q2. Is F1 = 0.98 a statevector artifact, or does it survive shot noise and depolarising noise?}
\textbf{Basis.} Paper uses Aer QasmSimulator (Sec.~IV.B), but the default FidelityQuantumKernel
sampler in qiskit-ml 0.9 is effectively statevector unless shots are set. A
quantum-kernel result that vanishes under 1000-shot sampling or 0.5\% two-qubit
depolarising error is not deployable.\\
\textbf{Next steps.} Rerun QSVC/Z on Aer with shots $\in \{512, 1024, 4096, 16384,
\infty\}$ and a FakeManila / FakeMumbai noise model at scale $\in \{0.5\times, 1\times,
2\times\}$. Plot F1 vs.\ shots and F1 vs.\ noise-scale. Free endpoint: local Aer.

\subsection*{Q3. Would ML-driven feature-map / feature-subset selection beat the hand-picked design?}
\textbf{Basis.} Paper picks features by PCA + heatmap and feature maps from three
canonical ones; both are ad-hoc. Kernel-target-alignment optimisation (Hubregtsen et
al.\ 2022, Glick et al.\ 2024) shifts QSVC F1 by 0.05--0.15 on comparable small
datasets.\\
\textbf{Next steps.} Sweep entangler patterns (linear / circular / full / pairwise) at
reps $\in \{1,2,3\}$, score by kernel-target alignment on the training split, evaluate
top-3 on test. Also do feature-subset ablations (6-choose-$k$ for $k \in \{3,4,5\}$)
to see whether the 4-feature pick actually dominates. Budget $\sim$30 min per ansatz.

\subsection*{Q4. Does QSVC/Z survive against modern quantum-kernel methods?}
\textbf{Basis.} Paper compares three fixed feature maps only. Since 2023, projected
quantum kernels (Huang et al.\ 2021) and trainable data-encoding kernels (Lloyd et al.\
2020) have become the honest state of the art for small-data quantum classifiers. A
REPLICATED result on a superseded kernel family is a data point but not a
competitiveness claim.\\
\textbf{Next steps.} Implement (a) a projected quantum kernel with projection dimension
$\in \{2,4,8\}$ on the same 200-record slice, (b) a trainable data-scaling variant of
ZFeatureMap where each rotation angle carries a learnable scalar, optimised jointly
with the SVM margin. Compare F1 to QSVC/Z. Free endpoint: local Aer or Sophia H100.

\subsection*{Q5. Does QML specialise usefully to fraud subtypes, or is the balanced-subset F1 hiding cliffs?}
\textbf{Basis.} Paper reports one aggregate F1 per cell. Real fraud detection needs
per-segment performance because losses are dominated by a few high-value fraud
subtypes. The balanced-subset design makes segment analysis impossible on the paper's
slice.\\
\textbf{Next steps.} On the full BankSim file, stratify by (a) merchant category, (b)
transaction-amount decile, (c) customer age bucket. Report QSVC/Z F1 and PR-AUC per
stratum vs.\ classical baseline. Identify subtypes where quantum wins by more than
$2\sigma$ of the seed distribution (20-seed sweep, $\sim$10 h CPU).
